\section{A sharp Euler constant at every positive real order}
\label{univeuler:sec:main}
The Gaussian maximum does not bound arbitrary Euler errors merely because
it bounds the value of the function. A contraction inequality for the
Euler kernel supplies the missing implication. It removes the divergent
parameter-dependent budgets near the signed-measure threshold, while
retaining the stronger constant one on the earlier domain $a\ge1$.

\begin{lemma}[A difference-of-powers bound]\label{univeuler:lem:powers}
For $m\ge2$, $0\le r\le1$ and $0\le v\le1$,
\[
 0\le(1-rv)^m-(1-r)^m\le\frac{1-v}{1+v}.
\]
For $0<v<1$, equality on the right is possible only when $m=2$,
and then it holds precisely at $r=1/(1+v)$.
\end{lemma}
\begin{proof}
The endpoints $v=0,1$ are immediate. Fix $0<v<1$ and put
$q=v^{1/(m-1)}$. Differentiation shows that the unique maximum over
$r\in[0,1]$ is at $r=(1-q)/(1-vq)$, where its value is
\[
 \frac{(1-v)^m}{(1-v^{m/(m-1)})^{m-1}}.
\]
Let $p=m/(m-1)\in(1,2]$. The function
$p\mapsto\log(1-v^p)$ has strictly negative second derivative
$-(\log v)^2v^p/(1-v^p)^2$. Its chord between $p=1$ and $p=2$
therefore gives
\[
 1-v^p\ge(1-v)(1+v)^{p-1},
\]
strictly when $1<p<2$. Substitution in the maximum proves the bound
and both equality assertions.
\end{proof}

\begin{lemma}[Contraction of the Euler tail kernel]
\label{univeuler:lem:contraction}
For an integer $N\ge1$, set $H_N(t)=(1-t)^N/(1+t)$ on $[0,1]$.
For $0\le r,v\le1$,
\begin{equation}\label{univeuler:eq:contraction}
 0\le H_N(rv)-H_N(r)\le\frac{1-v}{1+v}.
\end{equation}
If $0<v<1$, equality on the right holds precisely when $N=1$ and $r=1$.
\end{lemma}
\begin{proof}
For $N=1$, the difference between the right side and the middle
expression factors exactly as
\[
 \frac{(1-v)(1-r)(1-rv)}{(1+v)(1+rv)(1+r)}.
\]
For $N\ge2$, the uniformly convergent positive expansion
\[
 H_N(t)=\sum_{k\ge0}\frac{(1-t)^{N+k}}{2^{k+1}}
\]
expresses its difference as a probability-weighted sum of the preceding
power differences. Each is bounded as claimed, and the terms with
$N+k>2$ make the upper inequality strict for $0<v<1$. Decrease of
$H_N$ proves the lower inequality.
\end{proof}

Let $\omega_b$ be the probability measure
\[
 d\omega_b(v)=\frac{(-\log v)^{b-1}}{\Gamma(b)}\,dv,
 \qquad 0<v<1.
\]
Let $\pi_a$ be the law of $X=e^{-T_a}$ with
$T_a\sim\mathrm{Gamma}(a,1)$ when $a>0$, and put $\pi_0=\delta_1$.
These measures are independent. They retain ordinary probability mass
throughout the parameter domain, including the critical and subcritical
regions where the unweighted one-variable moment measure is not finite.

\begin{theorem}[All-order Euler representation and the sharp universal constant]
\label{univeuler:thm:sharp}
For every $a\ge0,b>0$, define $d_k$ and $E_N$ from the coefficients
$c_{2n+1}$ as in \eqref{rw:eq:eulerdef}, and use the analytic value
$g_{a,b}=\Imag F_{a,b}(i)$. Then $d_0=0$, $d_k<0$ for $k\ge1$,
and $E_N$ decreases strictly to $g_{a,b}$. Its exact error is
\begin{equation}\label{univeuler:eq:tail}
 2^N(E_N-g_{a,b})=
 \int\!\!\int
 \frac{H_N(x^2v^2)-H_N(x^2)}{1-v}\,d\pi_a(x)\,d\omega_b(v).
\end{equation}
Consequently, for every $N\ge1$,
\begin{equation}\label{univeuler:eq:axisbudget}
 0<2^N(E_N-g_{a,b})\le C(b),\qquad
 C(b)=\beta(b)+2^{-b}\eta(b).
\end{equation}
Equality in the upper bound holds exactly when $a=0$ and $N=1$.
The sharp uniform constant over all nonnegative outer orders is
\[
 C_* = C(b_*),
\]
where $b_*$ is the unique nondegenerate maximum from
Theorem~\ref{subpick:thm:axis-maximum}. It satisfies
\[
 \frac{227}{200}<C_*<\frac{1+\sqrt2}{2}<\frac54.
\]
For positive $a,b$ the strict inequality
$0<E_N-g_{a,b}<C_*2^{-N}$ holds, and $C_*$ is its least possible
uniform constant. When $a=0$ is included, the non-strict upper bound
is attained only at $a=0,b=b_*,N=1$.
\end{theorem}
\begin{proof}
The elementary moment laws give, for $n\ge0$,
\[
 c_{2n+1}=\int\!\!\int x^{2n}
             \frac{1-v^{2n}}{1-v}\,d\pi_a(x)\,d\omega_b(v).
\]
Indeed the two factors integrate to $(2n+1)^{-a}$ and $H_{2n}^{(b)}$.
The quotient is a finite geometric sum, so this identity includes $n=0$
without an endpoint singularity. Finite binomial expansion now gives
\[
 d_k=\int\!\!\int
       \frac{(1-x^2)^k-(1-x^2v^2)^k}{1-v}
                 \,d\pi_a(x)\,d\omega_b(v).
\]
For $k\ge1$ its integrand is strictly negative and has absolute value
at most $2k$, by the mean-value theorem. Finite summation is therefore
legitimate for every pair of parameters.

The Gaussian double kernel, or the axis formula when $a=0$, gives
\[
 g_{a,b}=\int\!\!\int
 \frac{(1+x^2)^{-1}-(1+x^2v^2)^{-1}}{1-v}
                    \,d\pi_a(x)\,d\omega_b(v).
\]
Its integrand is $-x^2(1+v)/((1+x^2)(1+x^2v^2))$, hence bounded.
Subtract the finite geometric sum defining $E_N$ from this identity
to obtain \eqref{univeuler:eq:tail}; no divergent boundary series is used.

Apply Lemma~\ref{univeuler:lem:contraction} with $r=x^2$ and its
variable $v$ replaced by $v^2$. The resulting majorant is
\[
 \frac{1-v^2}{(1-v)(1+v^2)}=\frac{1+v}{1+v^2}.
\]
Its $\omega_b$ integral is $C(b)$ by the axis identity.
It is bounded by $(1+\sqrt2)/2$, so the exact error tends to zero
at least geometrically. The strict signs of $d_k$ give the strict
decrease of $E_N$. The contraction equality criterion and the positive
densities of both probability measures prove the fixed-$b$ equality claim.

The unique axis maximum supplies the uniform constant and its bounds.
Sharpness over $a,b>0$ follows by fixing $b=b_*$, taking $a\downarrow0$
at $N=1$, and using bounded convergence in the Gaussian probability
integral. Thus the supremum is approached at the axis, while inclusion
of the axis gives the stated unique attaining triple.
\end{proof}

This theorem closes the previously open universal real-order Euler-budget
question. The separate contraction argument is essential: Gaussian
magnitude alone would not control later scaled errors. In particular,
this result does not say that their maximum at each positive parameter
pair is the first error, nor that the scaled errors always decrease.
The constant one remains sharper on $a\ge1$, and the sharp critical
constant $\pi/4+\log2/2$ remains sharper on $a+b\le1$.
The decimal diagnostic $C_*\approx1.13656110333950956095$ inherits the
axis experiment's status; the exact definition and inequalities above
are proved without relying on those digits.
